% Methodology for the final single-stage TIDE model.
% Required packages: amsmath, amssymb, bm

\section{Methodology}
\label{sec:methodology}

\subsection{Problem Formulation}

We consider causal estimation of gripper state, current hand position, screen
location, and short-horizon motion from four surface electromyography (sEMG)
channels and four six-axis inertial measurement units (IMUs). At time $t$, the
raw wearable observation is
\begin{equation}
    \mathbf{x}_t=[\mathbf{e}_t,\mathbf{a}_t,\boldsymbol{\omega}_t],
\end{equation}
where $\mathbf{e}_t\in\mathbb{R}^{4}$ contains the sEMG measurements and
$\mathbf{a}_t,\boldsymbol{\omega}_t\in\mathbb{R}^{12}$ contain acceleration and
angular velocity from the four IMUs. Using only observations available at or
before $t$, the model predicts the open/close state $\hat g_t$, current
Cartesian position $\hat{\mathbf p}_t\in\mathbb{R}^{3}$, normalized screen
location $\hat{\mathbf u}_t\in[0,1]^2$, and future Cartesian positions
$\hat{\mathbf p}_{t+h\mid t}$ for
$h\in\mathcal H_s=\{10,20,\ldots,200\}$ ms. Orientation, final endpoint,
spatial-cell, within-cell offset, and one-second Cartesian prediction heads are
not included in the final model.

The complete network is optimized end to end in a single stage. This choice was
made after a controlled comparison against separately pretrained and frozen
encoders: joint optimization preserved gripper-state performance while
improving current position, screen prediction, and future motion. The resulting
formulation allows state and motion supervision to shape a common causal
representation instead of fixing the feature space before the motion tasks are
introduced.

\subsection{Wearable Signals and Causal Preprocessing}
\label{subsec:preprocessing}

The four wearable units are distributed around the forearm. Each unit provides
one sEMG channel, three-axis acceleration, and three-axis angular velocity,
giving four sEMG and 24 IMU channels. Measurements are ordered by acquisition
timestamp and resampled to a uniform 100-Hz grid. A missing measurement is
forward-filled only when the most recent past observation is no more than
20 ms old. Longer gaps remain invalid and are represented explicitly by binary
validity indicators.

Each sEMG channel is processed with a fourth-order causal Butterworth band-pass
filter with lower cutoff 20 Hz and upper cutoff
\begin{equation}
    f_{\mathrm{high}}=\min(450,0.4f_s)\ \mathrm{Hz},
\end{equation}
where $f_s$ is the acquisition rate. Trailing root-mean-square envelopes of
20 and 50 ms are computed, producing eight sEMG activation features. The short
window preserves rapid recruitment changes, whereas the longer window provides
a stable activation estimate. Each IMU channel is processed by a trailing
three-sample median filter followed by a second-order causal 15-Hz low-pass
filter. The DC component is retained.

Means and standard deviations are estimated independently for every channel
using valid training samples only. Standardized values are clipped to
$[-20,20]$ and concatenated with their validity indicators, producing
\begin{equation}
    \tilde{\mathbf e}_t\in\mathbb{R}^{16},\qquad
    \tilde{\mathbf i}_t\in\mathbb{R}^{48}.
\end{equation}
A fused loss is evaluated only when at least 75\% of both sensor modalities are
valid. VIVE position is used as supervision for current and future Cartesian
position but is never provided as a model input.

\subsection{Causal Multiscale Encoders}
\label{subsec:encoders}

sEMG and IMU are encoded separately because they measure different physical
processes and exhibit different useful temporal scales. sEMG provides direct,
non-stationary evidence of muscle recruitment, whereas IMU supplies the
mechanical context needed to distinguish grasp-related activation from
reaching, rotation, and stabilization. Early concatenation would force both
signals through identical temporal filters before their modality-specific
structure had been extracted.

Each input stream is projected to width $d=128$. A frame-resolution path uses
causal depthwise convolutions with kernels of five and nine samples, covering
50 and 90 ms at the model rate. In parallel, a context path forms overlapping
16-frame patches with stride four and passes them through a four-layer,
four-head Transformer. Left padding ensures that the token associated with
time $t$ contains no later sample, and a causal attention mask blocks future
tokens. The most recent completed patch representation is held between patch
endpoints and combined with the local representation. The local path preserves
rapid state transitions, while the patch path models longer movement context.

Fusion is performed separately for the local and contextual paths. For
$r\in\{\mathrm{loc},\mathrm{ctx}\}$, let $\mathbf h_t^{e,r}$ and
$\mathbf h_t^{i,r}$ denote the sEMG and IMU features. A detached gate input
produces modality weights
\begin{equation}
 \boldsymbol{\alpha}_t^r=
 \operatorname{softmax}\!\left(
 g_r\!\left(\operatorname{sg}
 [\mathbf h_t^{e,r},\mathbf h_t^{i,r}]\right)\right),
\end{equation}
and the fused representation is
\begin{align}
 \mathbf z_t^r=\phi_r\big[&
 \alpha_t^{e,r}\mathbf h_t^{e,r},
 \alpha_t^{i,r}\mathbf h_t^{i,r},\nonumber\\
 &\mathbf h_t^{e,r}-\mathbf h_t^{i,r},
 \mathbf h_t^{e,r}\odot\mathbf h_t^{i,r}\big].
\end{align}
The weighted features preserve modality-specific evidence, while the difference
and product expose cross-modal disagreement and agreement. Stop-gradient is
applied only at the gate input; all task losses continue to train both modality
encoders through the fused representation.

\subsection{Gripper-State Prediction}
\label{subsec:state}

The state decoder combines a stable contextual prediction with a
high-resolution local prediction,
\begin{equation}
    \boldsymbol{\ell}_t^g=
    \boldsymbol{\ell}_t^{\mathrm{ctx}}
    +\sigma(\beta)\boldsymbol{\ell}_t^{\mathrm{loc}},
\end{equation}
where the local term is produced by an additional five-frame causal depthwise
convolution. The blend parameter is initialized with $\beta=-2$, initially
favoring context while allowing learning to sharpen open/close transitions.

A causal sEMG-only pathway supplies an additional correction to the state
logits. This pathway prevents the state estimate from depending exclusively on
mechanical motion and retains a direct route from muscle recruitment to the
open/close decision. Its output projection is initialized to zero, so it begins
as a no-op and is retained only when it improves the shared classification
objective. Learned MASK state tokens are used in this pathway during inference;
no ground-truth state is supplied to the network.

The posterior state probability
$\mathbf q_t=\operatorname{softmax}(\boldsymbol{\ell}_t^g)$ is used both as the
classification output and as a differentiable conditioning signal for the
motion pathway. Unlike the earlier frozen formulation, motion gradients may
therefore refine the encoder features and state-conditioned representation.

\subsection{State-Conditioned Residual Motion Model}
\label{subsec:motion}

The modality-specific contextual features
$\mathbf u_t^e,\mathbf u_t^i\in\mathbb{R}^{128}$ are passed to separate
two-layer causal GRUs. These recurrent layers model motion history while
retaining the multiscale features extracted by the patch encoders:
\begin{align}
 \mathbf h_{1:t}^{e}&=\mathbf u_{1:t}^{e}
 +W_a^e\operatorname{GRU}_e(\mathbf u_{1:t}^{e}),\\
 \mathbf h_{1:t}^{i}&=\mathbf u_{1:t}^{i}
 +W_a^i\operatorname{GRU}_i(\mathbf u_{1:t}^{i}).
\end{align}
The residual projections $W_a^e$ and $W_a^i$ are initialized to zero. Thus,
optimization begins from the causal encoder features and gradually introduces
recurrent corrections rather than replacing the representation at
initialization.

The state posterior is embedded as $\mathbf s_t=\phi_s(\mathbf q_t)$. IMU
provides a mechanical base, and the sEMG stream supplies a bounded,
state-conditioned correction:
\begin{align}
 \mathbf c_t&=\tanh\!\left(\phi_c
 [\mathbf h_t^e,\mathbf s_t]\right),\\
 \gamma_t&=\sigma\!\left(\phi_\gamma
 [\mathbf h_t^i,\mathbf h_t^e,\mathbf s_t]\right),\\
 \mathbf f_t&=\operatorname{LN}
 (\mathbf h_t^i+\gamma_t\mathbf c_t).
\end{align}
The correction projection is zero-initialized and the gate bias is initialized
to $-2$. This gives the network a conservative IMU-centered starting point
while allowing sEMG to modify the motion representation when supported by the
joint objective.

\subsection{Prediction and Auxiliary Heads}

Two-layer multilayer perceptrons map $\mathbf f_t$ directly to current
Cartesian position and normalized screen location. Screen coordinates are
converted using the horizontal and vertical display dimensions independently,
\begin{equation}
    \hat{\mathbf u}_t^{\mathrm{px}}
    =(W\hat u_{x,t},H\hat u_{y,t}),
\end{equation}
with $W=1920$ and $H=1080$ in our experiments. No spatial-grid classification
or within-cell offset head is used.

For future prediction, the decoder concatenates $\mathbf f_t$ with the horizon
encoding
\begin{equation}
 \mathbf b(\tau)=[\tau,\tau^2,\sin(\pi\tau),\cos(\pi\tau)]
\end{equation}
and predicts absolute Cartesian position at every 10-ms step through 200 ms.

A separate training-only decoder receives the sEMG representation and a
horizon encoding at 100-ms intervals through 1 s. It predicts future IMU
changes and is omitted from the deployed output interface. This auxiliary task
encourages the sEMG pathway to preserve neuromuscular activity that anticipates
subsequent mechanical motion; the final model does not output a one-second
Cartesian trajectory.

\subsection{Unified Learning Objective}
\label{subsec:loss}

Let $\rho(\cdot)$ denote element-wise smooth-$L_1$ and let
$\langle\cdot\rangle_{\mathcal V}$ average only valid targets. The gripper loss
is class-balanced cross entropy,
\begin{equation}
 \mathcal L_g=-\frac{1}{|\mathcal V_g|}
 \sum_{t\in\mathcal V_g}w_{g_t}
 \log\operatorname{softmax}(\boldsymbol{\ell}_t^g)_{g_t},
 \qquad w_c=\frac{N}{2N_c}.
\end{equation}
Current position is supervised in the training-standardized Cartesian space:
\begin{equation}
 \mathcal L_{\mathrm{pos}}=
 \left\langle\rho(\hat{\mathbf p}_t-\mathbf p_t)\right\rangle_{\mathcal V_p}.
\end{equation}

For the direct screen head, normalized residuals are scaled independently by
display width and height. Defining
$\mathbf D=\operatorname{diag}(W/100,H/100)$ and normalized trial progress
$r_t\in[0,1]$, the screen objective is
\begin{equation}
 \mathcal L_{\mathrm{pixel}}=
 \frac{\sum_{t\in\mathcal V_u}a_t
 \rho\!\left(\mathbf D
 (\hat{\mathbf u}_t-\mathbf u_t)\right)}
 {\sum_{t\in\mathcal V_u}a_t},
 \qquad a_t=0.25+3.75r_t^4.
\end{equation}
The separate $x$ and $y$ scaling respects the 16:9 display geometry, while
$a_t$ emphasizes the later part of the reach, when the intended screen target
should have converged.

Dense future-position supervision is
\begin{equation}
 \mathcal L_{\mathrm{future}}=
 \frac{1}{|\mathcal H_s|}
 \sum_{h\in\mathcal H_s}
 \left\langle
 \rho(\hat{\mathbf p}_{t+h\mid t}-\mathbf p_{t+h})
 \right\rangle_{\mathcal V_{t,h}}.
\end{equation}
This trains the trajectory directly instead of deriving it from the current
position or a constant-velocity assumption.

The only long-horizon auxiliary objective is sEMG-to-future-IMU prediction. Let
$\mathcal H_l=\{100,200,\ldots,1000\}$ ms, let $h^{-}$ denote the preceding
horizon, and let $\Delta\bar{\mathbf i}_{t,h}$ be the mean IMU value over the
interval $(t+h^{-},t+h]$ minus the current IMU value. Then
\begin{align}
 \mathcal L_{e\rightarrow i}^{(h)}
 &=\left\langle\rho\!\left(
 \widehat{\Delta\bar{\mathbf i}}_{t,h}
 -\Delta\bar{\mathbf i}_{t,h}\right)\right\rangle,\\
 \mathcal L_{e\rightarrow i}
 &=\frac{\sum_{h\in\mathcal H_l}\lambda_h
 \mathcal L_{e\rightarrow i}^{(h)}}
 {\sum_{h\in\mathcal H_l}\lambda_h},\\
 \lambda_h&=\exp\!\left(-\frac{h}{500\ \mathrm{ms}}\right).
\end{align}
The auxiliary decoder receives no IMU representation. It therefore cannot
solve the task by copying inertial features and must instead extract sEMG
patterns associated with later mechanics. Future IMU values are used only to
construct training targets.

The complete end-to-end objective is
\begin{equation}
 \boxed{\begin{aligned}
 \mathcal L={}&\mathcal L_g+\mathcal L_{\mathrm{pos}}
 +0.35\mathcal L_{\mathrm{pixel}}\\
 &+0.50\mathcal L_{\mathrm{future}}
 +0.05\mathcal L_{e\rightarrow i}
 \end{aligned}}
\end{equation}
No orientation, endpoint, grid, offset, arrival-consistency, correction
regularization, masked-sEMG reconstruction, or one-second Cartesian intent
loss is used.

\subsection{Optimization and Causal Inference}

The network is trained end to end with AdamW using learning rate
$3\times10^{-4}$, weight decay $10^{-4}$, batch size four, and gradient-norm
clipping at one. Training runs for at most 60 epochs and stops after 12
validation epochs without improvement. Checkpoints minimize the validation
score
\begin{equation}
 S_{\mathrm{val}}=E_{\mathrm{pos}}^{\mathrm{cm}}
 +0.01E_{\mathrm{pixel}}^{\mathrm{px}}
 +5(1-F_1^{\mathrm{state}}).
\end{equation}
Trials, rather than overlapping frames, are assigned to 60\% training, 20\%
validation, and 20\% test partitions using a fixed split seed. All
normalization statistics and class weights are derived from the training
partition. The selected model contains 3,242,105 trainable parameters.

At inference, the same causal preprocessing is applied with a two-second
history and a 300-ms warm-up. State output uses hysteresis: the gripper changes
to closed when $P(g_t=\mathrm{close})\geq0.7$, changes to open when
$P(g_t=\mathrm{close})\leq0.3$, and otherwise retains its previous state. Each
valid output contains state probabilities, current Cartesian position,
continuous screen location, and 20 future Cartesian positions through 200 ms.
No VIVE sample, future wearable measurement, target label, or online parameter
update is required.
